% TOP_PROBLEM: {"id": 3230, "title": "Coloring the Odd Distance Graph", "category_id": 3, "category": {"id": 3, "name": "graph_theory", "display_name": "Graph Theory", "description": "Problems involving graphs, networks, and their properties.", "slug": "graph-theory", "order_index": 3, "created_at": "2026-07-31T15:26:25.671Z"}, "status": "open"}
% FIRST_POSTED: null
% ENTRY_KIND: statement_only
% Imported from: batch02.tex; UnsolvedMath ID 3230
% Compile twice: lualatex 3230.tex
\documentclass[11pt]{article}
\usepackage{fontspec}
\setmainfont{Latin Modern Roman}
\usepackage{amsmath,amssymb,mathtools,unicode-math}
\setmathfont{Latin Modern Math}
\usepackage{xurl,hyperref}
\hypersetup{hidelinks,pdftitle={UnsolvedMath Problem 3230}}
\newcommand{\sref}[2]{\hyperlink{source:#1:#2}{[S#2]}}
\newcommand{\cataloguescope}[1]{\paragraph{Mathematical question.}#1}
\begin{document}
% UnsolvedMath ID 3230; source catalogue code OPG-616
\setcounter{section}{3229}
\section{Coloring the Odd Distance Graph}\label{3230}
{\small\sffamily UnsolvedMath ID 3230 · Graph Theory}
\subsection{Definitions and mathematical statement}

\paragraph{Shared notation.}$\mathbb N=\{1,2,\ldots\}$, and $\mathbb Z,\mathbb Q,\mathbb R,\mathbb C$ denote the integers, rationals, reals and complex numbers. Any other conventions are specified locally.


The odd-distance graph $\mathcal O$ has vertex set $\mathbb R^2$ and joins distinct points $x,y$ when their Euclidean distance is an odd positive integer. A proper coloring assigns different colors to adjacent points; no measurability restriction is imposed. A measurable finite coloring means that each color class is Lebesgue measurable.
\paragraph{Statement recorded in the supplied source.}
Is there no proper coloring $\mathbb R^2\to\{1,\ldots,m\}$ of $\mathcal O$ for any finite $m$, that is, is $\chi(\mathcal O)$ infinite? The accompanying measurable-coloring variant asks the same question with measurable color classes. \sref{3230}{2}
\subsection{Short English statement}
Must every finite coloring of the plane contain two equally colored points at an odd integer distance, even when color classes are arbitrary sets?
\subsection{Sources}
\begin{description}
\item[\hypertarget{source:3230:1}{S1}] ULAM.AI and the UnsolvedMath Contributors, UnsolvedMath: A Curated Collection of Open Mathematics Problems, record 3230 (OPG-616). CC BY 4.0. \url{https://huggingface.co/datasets/ulamai/UnsolvedMath}.
\item[\hypertarget{source:3230:2}{S2}] Open Problem Garden, Coloring the Odd Distance Graph, original node 616, statement and mathematical discussion. \url{https://www.openproblemgarden.org/op/coloring_the_odd_distance_graph}.
\end{description}
\label{end:3230}
\end{document}
